
This study compared Neural Functional Transformer (NFT) encoders with flat-MLP baselines for generalization gap prediction on the Unterthiner model zoo under permutation stress. NFT exhibited near-zero permutation sensitivity (delta-rho approximately 4.7 x $10^-7$) across 21 training runs, while flat-MLP showed significant degradation (delta-rho = 0.160--0.640). Mediation analysis (delta-R-squared = 0.228) indicated that the robustness difference was associated with equivariant attention capturing structural signals beyond what augmentation provides. Permutation augmentation partially reduced flat-MLP sensitivity but with high seed-dependent variance. L2-norm canonicalization collapsed to constant predictions.

These results were obtained on a single model zoo (Unterthiner MNIST CNN zoo adapted to token format, 29,997 models) with a single prediction target (generalization gap). Two of five planned sub-hypotheses (cross-pipeline transfer and graceful degradation) were not executed. A matched-parameter comparison between NFT and flat-MLP was not conducted.

Future work includes: (1) executing the cross-pipeline transfer experiments (h-m3, h-m4) using existing infrastructure, (2) validating on native FC-MLP zoo weights, (3) testing stronger canonicalization approaches (Hungarian alignment, sort-by-magnitude), (4) conducting a matched-parameter comparison to disentangle equivariance effects from parameter count effects, and (5) extending to additional model zoos and prediction targets.
